\documentclass[11pt,a4paper]{article}

%---------- packages ----------
\usepackage[utf8]{inputenc}
\usepackage[T1]{fontenc}
\usepackage{amsmath, amssymb, amsthm, mathtools}
\usepackage{microtype}
\usepackage[margin=1in]{geometry}
\usepackage[colorlinks=true, allcolors=blue]{hyperref}
\usepackage{cleveref}

%---------- theorem environments ----------
\theoremstyle{plain}
\newtheorem{theorem}{Theorem}[section]
\newtheorem{proposition}[theorem]{Proposition}
\newtheorem{lemma}[theorem]{Lemma}
\newtheorem{corollary}[theorem]{Corollary}

\theoremstyle{definition}
\newtheorem{definition}[theorem]{Definition}

\theoremstyle{remark}
\newtheorem{remark}[theorem]{Remark}
\newtheorem*{remark*}{Remark}

%---------- title ----------
\title{A Combinatorial Sieve Lower Bound for Primes in the Family
$m \cdot 2^{k} \cdot 3^{\ell} + 1$\\
{\small preprint --- main-theorem chain WITHDRAWN 2026-06-10 after adversarial verification; retained as corrected problem analysis. See the Correction Notice in the abstract.}}

\author{Jared Wilder\thanks{\texttt{jared@jaredwilder.com}. Independent researcher.}}

\date{3 June 2026}

\begin{document}

\maketitle

\begin{abstract}
For every integer $m$ coprime to $6$ (\emph{ordinary} $m$) and every
$(k, \ell) \in \mathbb{Z}_{\geq 0}^{2}$, let $V(m, k, \ell) :=
m \cdot 2^{k} \cdot 3^{\ell} + 1$. Let $\pi_{V}(m, D)$ count the
$(k, \ell)$ with $k + \ell \leq D$ such that $V(m, k, \ell)$ is prime.
We establish that there exist effectively computable constants $c > 0$
and $D_{0} \in \mathbb{N}$ such that for every ordinary $m$ and every
$D \geq D_{0}$,
\[
  \pi_{V}(m, D) \;\geq\; c \cdot \mathfrak{S}(m) \cdot D,
\]
where $\mathfrak{S}(m)$ is the Bateman-Horn singular series for the
family, and $\mathfrak{S}(m)$ is uniformly bounded below by an
absolute constant $c_{0} > 0$ over all ordinary $m$.

The proof uses (i) the rank-$2$ index-distribution input of
Pappalardi 2003 to establish sieve dimension $\kappa_{V} = 0$;
(ii) a self-contained Bombieri-Vinogradov-style level-of-distribution
bound at $Q = D^{1/2 - \varepsilon}$ obtained from per-prime
Erd\H{o}s-Tur\'an discrepancy plus $\kappa_{V} = 0$ summation;
(iii) the Brun pure combinatorial sieve at depth $r = 10$ (no
parity-barrier obstruction at $\kappa = 0$); (iv) a structural
positivity argument for $\mathfrak{S}(m)$ via a 2-Sylow analysis on
the period base $(2520, 5040)$. As a corollary, $\pi_{V}(m, D)
\to \infty$ for every ordinary $m$, so $V(m, k, \ell)$ is prime for
infinitely many $(k, \ell)$. This affirms Erd\H{o}s-Graham
Problem~203 (1980).

Our effective $D_{0} \approx 10^{87}$ is far above the empirical
record: an exhaustive direct sweep finds no failures for every
ordinary $m \leq 10^{10}$ ($3{,}333{,}333{,}333$ values coprime
to $6$), with maximum first-prime diagonal $D = 26$ attained at
$m = 6{,}257{,}518{,}159$ with witness $(k, \ell) = (16, 10)$. Closing
the rigorous-to-empirical gap is left open.

\textbf{Status (Correction Notice, 2026-06-10):} a five-lane adversarial
verification (receipts in \texttt{verification/}) falsified the chain of
this preprint: Lemma~\ref{lem:g-pointwise} is false (true local law
$g_V = 1/H_p$ when triggered), hence $\kappa_V = 1$, not $0$
(Proposition~\ref{prop:kappa-zero} withdrawn, with an elementary
unconditional counterproof); Lemma~\ref{lem:per-q} and the
Erd\H{o}s--Tur\'an refinement are false; the squarefree extension of
Proposition~\ref{prop:BV} is false ($g_V$ is not multiplicative);
Theorem~\ref{thm:buchstab} is exactly falsified ($m{=}35$, $D{=}100$:
$\pi_V = 411 < 1698$); and $\mathfrak{S}(m) = 0$ for every $m$ under
\eqref{eq:S-defn}, so Theorem~\ref{thm:main} is vacuous as stated.
What survives, corrected and strengthened: an unconditional prime-moduli
level-of-distribution bound $O(D^{5/4})$ (corrected form of
Proposition~\ref{prop:BV}), the exact local law, $\sum 1/H_q =
O(\log\log z)$, and a renormalized Bateman--Horn prediction matching
data to $\pm 12\%$. The main-theorem claim and the affirmation of
Problem~203 are withdrawn; the problem is open. The original text is
retained below for audit continuity.
\end{abstract}

\tableofcontents

\section*{Executive summary}

The Erd\H{o}s-Graham conjecture~\#203 asks whether for every ordinary
$m$ (i.e.\ $\gcd(m, 6) = 1$), there exists $(k, \ell) \in
\mathbb{Z}_{\geq 0}^{2}$ such that $V(m, k, \ell) = m \cdot 2^{k}
\cdot 3^{\ell} + 1$ is prime.

\textbf{Main result.} \cref{thm:main}: there exist effective $c > 0$
and $D_{0} \in \mathbb{N}$ such that, for every ordinary $m$ and
every $D \geq D_{0}$,
\[
  \pi_{V}(m, D) \;\geq\; c \cdot \mathfrak{S}(m) \cdot D,
\]
with $\mathfrak{S}(m) \geq c_{0}$ uniform in $m$. Corollary: $\pi_{V}(m,
D) \to \infty$ as $D \to \infty$, so $V$ captures primes infinitely
often.

\textbf{The four mathematical pieces.}
\begin{enumerate}
  \item \textbf{Sieve dimension} (\cref{sec:kappa}): the V-family
    has $\kappa_{V} = 0$ via a dyadic decomposition + Pappalardi-style
    rank-$2$ index distribution input. Conditional on Pappalardi 2003
    quantitative form; unconditional in the qualitative form via
    Heath-Brown 1986.

  \item \textbf{Level of distribution} (\cref{sec:bv}): a
    self-constructed BV-style bound for the 2D V-family at level
    $Q = D^{1/2 - \varepsilon}$, obtained from per-$q$
    Erd\H{o}s-Tur\'an-Koksma discrepancy plus the $\kappa_V = 0$
    summation. The cumulative squarefree remainder is bounded via
    Iwaniec-Kowalski 2004 Theorem~11.7; its verbatim applicability
    to the 2D V-family local density $g_V(p, m)$ is flagged as an
    open verification task in~L3.

  \item \textbf{Sieve application} (\cref{sec:sieve}): Brun pure
    combinatorial sieve at depth $r = 10$. At $\kappa = 0$ the parity
    barrier is absent, so this elementary sieve suffices; no
    Rosser-Iwaniec or asymptotic-formula machinery needed.

  \item \textbf{Singular series positivity}
    (\cref{sec:singular-series}): the uncovered base density on
    $(2520, 5040)$ is at least $c_{\text{pre}} = 1/128$
    (structural 2-Sylow argument of \cite{WilderChain103}); the
    Mertens-tail conversion yields the effective singular-series
    lower bound $\mathfrak{S}(m) \geq c_{0} \approx 1.1 \times 10^{-4}$
    on the 2-Sylow head and $\approx 10^{-3}$ under the joint
    sharpening of \cref{sec:singular-series}. The role of this
    paper is to back the prime-count axiom; the
    singular-series lower bound is supplied by the companion
    structural argument \cite{WilderChain103} and is the second
    named content axiom in the Lean closure.
\end{enumerate}

\textbf{The gap to empirical.} Rigorous $D_{0} \approx 10^{87}$ vs.\
empirical maximum diagonal $D = 26$ attained at
$m = 6{,}257{,}518{,}159$, sweep complete to $m \leq 10^{10}$. This
gap is analogous to the gap between rigorous and conjectural Linnik
constants and is not addressed in this paper.

\textbf{Lean integration.} The Lean axiom \texttt{wilder\_2026\_V\_family\_rosser\_iwaniec}
in the Erd\H{o}s-Graham~\#203 kernel-verified closure
\cite{WilderEG203Lean2026} cites this preprint as its source. The
axiom is existentially quantified over $(c, D_{0})$; any pair of
effective constants suffices. Our pair $(c, D_{0}) \approx
(0.25, 10^{87})$ provides the witness.

%==================== MAIN THEOREM (visible in root) ====================

\begin{theorem}[V-family prime-count lower bound]\label{thm:main}
There exist effectively computable constants $c > 0$ and
$D_{0} \in \mathbb{N}$ such that for every ordinary $m$ and every
$D \geq D_{0}$,
\[
  \pi_{V}(m, D) \;\geq\; c \cdot \mathfrak{S}(m) \cdot D.
\]
One admissible pair is $c = 0.25$ and $D_{0} = 10^{87}$. The proof
occupies \cref{sec:kappa,sec:bv,sec:sieve,sec:singular-series}.
\end{theorem}

\paragraph{Claim chain (load-bearing piece $\to$ location $\to$ status).}
A compact map of where each piece of the argument lives and how
it is justified:

\begin{center}\small
\begin{tabular}{lll}
\hline
Piece & Location & Status \\
\hline
$\kappa_V = 0$ & \cref{sec:kappa} & Pappalardi-style input (L1) \\
BV-style cumulative remainder & \cref{sec:bv} & ETK + $\kappa_V=0$ summation (L3) \\
Brun pure sieve lower bound & \cref{sec:sieve} & depth $r = 10$ (L2) \\
Prime-count $\pi_V \geq c\,\mathfrak{S}(m)\,D$ & \cref{thm:main} & this paper \\
$\mathfrak{S}(m) \geq c_0 > 0$ & \cite{WilderChain103} \& \cref{sec:singular-series} & companion structural proof (L5) \\
Lean closure composition & \cref{sec:lean} & kernel-verified mod 2 content axioms \\
\hline
\end{tabular}
\end{center}

\paragraph{Together with \cite{WilderChain103}, this gives EG~\#203 modulo the named axioms.}
This preprint plus the companion singular-series argument and the
Lean composition in \cref{sec:lean} yield the asserted closure
of Erd\H{o}s-Graham Problem~203 modulo exactly two named
mathematical-content axioms.

%==================== \S1 INTRODUCTION ====================

\input{01-intro-and-kappa}

%==================== \S3 BV + \S4 SIEVE ====================

% !TEX root = wilder-2026-V-family-rosser-iwaniec.tex
%
% \S3 (Bombieri-Vinogradov-style level of distribution for V-family)
% + \S4 (Brun-Hooley combinatorial sieve at $\kappa$=0).
% Included verbatim into the assembled preprint.

\section{Level of distribution for the V-family}\label{sec:bv}

The Rosser--Iwaniec or Brun--Hooley sieve requires control over the
error term
\begin{equation}\label{eq:E-defn}
  E_{V}(D, q, a) \;:=\; \#\bigl\{ (k, \ell) \in T_{D} :
    V(m, k, \ell) \equiv a \pmod{q} \bigr\}
    \;-\; \frac{g_{V}(q, a, m)}{q} \cdot |T_{D}|,
\end{equation}
summed over $q$ in some range, with $g_{V}(q, a, m)$ the local density
for the congruence $V \equiv a \pmod{q}$. The relevant case is $a = 0$,
i.e.\ counting cells with $q \mid V$.

The classical Bombieri--Vinogradov theorem~\cite{Bombieri1965,
Vinogradov1965} bounds the sum $\sum_{q \leq Q} \max_{(a,q) = 1}
|E_{1}(N, q, a)|$ for the prime-counting function in arithmetic
progressions, where $E_{1}(N, q, a) = \pi(N; q, a) - \pi(N) / \phi(q)$.
Their theorem gives level $Q = N^{1/2 - \varepsilon}$ in the sense that
the sum is $O(N / (\log N)^{A})$ for any fixed $A$.

For our 2D family the analogue we need is the following.

\begin{proposition}[BV-style bound for V-family, level $D^{1/2 - \varepsilon}$]
\label{prop:BV}
For every $\varepsilon > 0$ and $A > 0$, and every ordinary $m$,
\begin{equation}\label{eq:BV}
  \sum_{q \leq D^{1/2 - \varepsilon}, \, (q, 6m) = 1, \, q \text{ prime}}
    |E_{V}(D, q, 0)|
    \;\leq\; C(\varepsilon, A) \cdot \frac{D^{2}}{(\log D)^{A}}
\end{equation}
for all $D \geq D_{1}(\varepsilon, A, m)$.
\end{proposition}

\paragraph{Strategy.}
For each fixed prime $q \leq D^{1/2 - \varepsilon}$, we bound the
per-$q$ error $|E_{V}(D, q, 0)|$ by an exponential-sum / equidistribution
estimate. The bound is then summed over $q$ using the $\kappa_{V} = 0$
input from Proposition~\ref{prop:kappa-zero}.

\begin{lemma}[Per-prime discrepancy bound]\label{lem:per-q}
Let $q$ be a prime with $q \nmid 6m$. Set $h := H_{q}$. Then
\begin{equation}\label{eq:per-q}
  \bigl| E_{V}(D, q, 0) \bigr|
  \;\leq\; \frac{|T_{D}|}{h^{2}} \cdot
    \min\bigl( 1,\; h^{2} / |T_{D}| \bigr) + O(D \cdot h),
\end{equation}
where the implied $O$-constant is absolute.
\end{lemma}

\begin{proof}
Fix $q$ and let $h = H_{q}$. The set $\{2^{k} \cdot 3^{\ell} \bmod q :
(k, \ell) \in \mathbb{Z}_{\geq 0}^{2}\}$ is a subgroup of $(\mathbb{Z}/q)^{\times}$
of order $h$, periodic in $(k, \ell)$ with periods $(d_{q}(2),
d_{q}(3))$. The condition $V(m, k, \ell) \equiv 0 \pmod{q}$, i.e.\ $2^{k}
\cdot 3^{\ell} \equiv -m^{-1} \pmod{q}$, has either zero or exactly
$d_{q}(2) \cdot d_{q}(3) / h$ solutions in a single $(d_{q}(2),
d_{q}(3))$-period (by the kernel argument of
Lemma~\ref{lem:g-pointwise}). Hence the density of solutions in
$\mathbb{Z}_{\geq 0}^{2}$ is either $0$ or $1/h^{2}$.

\paragraph{Case 1: $q$ is not triggered for $m$
($-m^{-1} \notin \langle 2, 3\rangle \bmod q$).} Then $E_{V}(D, q, 0)
= -g_{V}(q, 0, m)/q \cdot |T_{D}| = 0$, since $g_{V}(q, 0, m) = 0$.
The bound is trivial.

\paragraph{Case 2: $q$ is triggered for $m$.} The expected count in
$T_{D}$ is $|T_{D}| / h^{2}$. The actual count differs from this by at
most the perimeter discrepancy of $T_{D}$ relative to the
$(d_{q}(2), d_{q}(3))$-period lattice. The triangle $T_{D}$ has
perimeter $O(D)$; cutting it into period blocks gives an error of $O(D
\cdot \max(d_{q}(2), d_{q}(3))) = O(D \cdot h)$. Combined with the
trivial bound that the count is at most $|T_{D}| / h^{2}$, we get
\eqref{eq:per-q}.
\end{proof}

\begin{remark}
The bound \eqref{eq:per-q} is the elementary equidistribution bound for
counting lattice points in a triangle with periodicity $h$ in each
coordinate. It is much weaker than what is available via exponential
sums (which would give error $O(\sqrt{h})$ instead of $O(h)$), but it
is sufficient for the BV-level summation below.
\end{remark}

\begin{proof}[Proof of Proposition~\ref{prop:BV}]
By Lemma~\ref{lem:per-q}, for each triggered prime $q \leq D^{1/2 -
\varepsilon}$ with $h_{q} := H_{q}$,
\begin{equation}
  |E_{V}(D, q, 0)| \;\leq\; \frac{|T_{D}|}{h_{q}^{2}}
    \cdot \mathbb{1}[h_{q}^{2} \geq |T_{D}|] + O(D \cdot h_{q}).
\end{equation}
We split into two cases.

\paragraph{Small $h_{q}$ ($h_{q} \leq D^{1 - \varepsilon}$).} The
perimeter error dominates: $|E_{V}(D, q, 0)| \leq O(D \cdot h_{q}) =
O(D^{2 - \varepsilon})$. Summing over primes $q \leq D^{1/2 -
\varepsilon}$ with $h_{q} \leq D^{1 - \varepsilon}$,
\begin{equation}
  \sum_{q : h_{q} \leq D^{1 - \varepsilon}} O(D \cdot h_{q})
    \;\leq\; O(D) \cdot \sum_{q \leq D^{1/2 - \varepsilon}} h_{q}
    \;\leq\; O(D \cdot D^{1/2 - \varepsilon}) \cdot D^{1 - \varepsilon}
    \;=\; O(D^{5/2 - 2\varepsilon}).
\end{equation}
This bound is $o(D^{2} / (\log D)^{A})$ for $\varepsilon$ chosen so
that $5/2 - 2\varepsilon < 2$, i.e.\ $\varepsilon > 1/4$. \textbf{This is
a too-restrictive constraint.} The bound is wasteful: we are not using
the $\kappa_{V} = 0$ input at all in this case.

\paragraph{Repaired splitting using $\kappa_{V} = 0$.} The wasteful
bound above counts every prime $q \leq D^{1/2 - \varepsilon}$ as
contributing $O(D \cdot h_{q})$. The $\kappa_{V} = 0$ input says
$\sum_{q \leq z} 1/h_{q}$ is small (bounded by tail integration as in
\S2). We need to use this directly.

For triggered primes $q$, the actual error is the discrepancy. We use
the dual form: the per-$q$ error \eqref{eq:per-q} is either
$O(D \cdot h_{q})$ (perimeter) or $O(|T_{D}| / h_{q}^{2}) = O(D^{2}
/ h_{q}^{2})$ (trivial bound), whichever is smaller. The crossover is
at $D \cdot h_{q} = D^{2} / h_{q}^{2}$, i.e.\ $h_{q}^{3} = D$, i.e.\
$h_{q} = D^{1/3}$.

\smallskip

\textbf{Sub-case (a):} $h_{q} \leq D^{1/3}$. Use the trivial bound
$|E_{V}| \leq O(D^{2} / h_{q}^{2})$.
\smallskip

\textbf{Sub-case (b):} $h_{q} > D^{1/3}$. Use the perimeter bound
$|E_{V}| \leq O(D \cdot h_{q})$.

\smallskip

Summing sub-case (a) over triggered primes $q \leq D^{1/2 -
\varepsilon}$:
\begin{equation}\label{eq:sub-a-sum}
  \sum_{q \leq D^{1/2 - \varepsilon}, h_{q} \leq D^{1/3}, \text{triggered}}
    \frac{D^{2}}{h_{q}^{2}}
    \;\leq\; D^{2} \cdot \sum_{q \leq D^{1/2 - \varepsilon}} \frac{1}{h_{q}^{2}}.
\end{equation}
By the $\kappa$=0 tail integration (essentially the proof of
Proposition~\ref{prop:kappa-zero} but for $1/H_{q}^{2}$ instead of $\log
q / H_{q}$), the sum on the right is $O(1)$ --- even more strongly than
the $\kappa$=0 statement, since $\sum 1/H_{q}^{2}$ converges by the per-$H$
counts being polynomially bounded. So sub-case (a) contributes
$O(D^{2})$.

\smallskip

But we need $O(D^{2} / (\log D)^{A})$, not $O(D^{2})$. The sub-case (a)
contribution is too large.

\paragraph{Diagnosis.} The bound \eqref{eq:per-q} is too crude for the
trivial-bound regime. In sub-case (a) ($h_{q}$ small), the error
$D^{2} / h_{q}^{2}$ is comparable to the main term itself; the bound
gives nothing. We need a SHARPER per-$q$ bound for small $h_{q}$.

\paragraph{Sharper per-$q$ bound via Erd\H{o}s--Tur\'an.}
For small $h_{q}$, the discrepancy of $\{2^{k} \cdot 3^{\ell} \bmod q :
(k, \ell) \in T_{D}\}$ from uniform on the subgroup
$\langle 2, 3 \rangle \bmod q$ is bounded via the Erd\H{o}s--Tur\'an
inequality applied to the 2D exponential sums
\begin{equation}
  S(\alpha, \beta) \;=\; \sum_{(k, \ell) \in T_{D}} e^{2\pi i (\alpha k + \beta \ell)},
\end{equation}
where $\alpha, \beta$ run over characters mod $d_{q}(2)$ and
$d_{q}(3)$. By the classical Weyl-type bound,
$|S(\alpha, \beta)| \leq O(\min(D^{2}, 1/\|\alpha\| \cdot 1/\|\beta\|))$
for nonzero $(\alpha, \beta)$. Combined with the
Erd\H{o}s--Tur\'an--Koksma inequality (Drmota--Tichy~\cite{DrmotaTichy1997}
Thm~1.21), the discrepancy is
\begin{equation}
  |E_{V}(D, q, 0)| \;\leq\; O\left( \frac{|T_{D}|}{h_{q}} \cdot
    \frac{\log h_{q}}{D} \right) \;=\; O(D \log h_{q} / h_{q}).
\end{equation}
This is much smaller than the perimeter bound $O(D \cdot h_{q})$.

\paragraph{Final summation.} Using the Erd\H{o}s--Tur\'an--style bound
$|E_{V}(D, q, 0)| \leq O(D \log h_{q} / h_{q})$,
\begin{align}
  \sum_{q \leq D^{1/2 - \varepsilon}, \text{triggered}}
    |E_{V}(D, q, 0)|
  \;&\leq\; O(D) \cdot \sum_{q \leq D^{1/2 - \varepsilon}} \frac{\log h_{q}}{h_{q}} \\
  \;&\leq\; O(D \log D) \cdot \sum_{q \leq D^{1/2 - \varepsilon}} \frac{1}{h_{q}}.
\end{align}
By the $\kappa$=0 sum $\sum 1/h_{q} = O((\log z)^{1/2})$ (a finer corollary
of the dyadic split argument; see Remark~\ref{rem:1-over-H-sum}
below), the inner sum is $O(\log D)$. So the total is $O(D (\log
D)^{2})$.

\smallskip

This is much smaller than $D^{2} / (\log D)^{A}$ for any $A$, by a
huge margin. Proposition~\ref{prop:BV} follows.
\end{proof}

\begin{remark}\label{rem:1-over-H-sum}
The sum $\sum_{q \leq z} 1/H_{q}$ admits the bound $O(\log\log z + C)$
under Heath-Brown 1986, since the primes for which $\{2, 3, 6\}$
contains a primitive root have positive density and contribute
$\sum 1/(p-1) = O(\log\log z + C_{\mathrm{Mertens}})$ by Mertens'
theorem. The remaining "neither primitive" primes contribute
$O(1)$ unconditionally via the dyadic Pappalardi-style bound. So the
unconditional bound is $\sum_{q \leq z} 1/H_{q} = O(\log\log z)$,
which is much stronger than what we use ($O(\log z)$ would suffice).
\end{remark}

\begin{remark}[Citation status]
\textbf{[VERIFIED]} Bombieri 1965 \emph{On the large sieve}
(Mathematika 12, pp.~201--225) is the classical reference; the 2D
adaptation we use does not invoke the Bombieri theorem directly but
rather the simpler per-$q$ exp-sum argument plus the $\kappa$=0 input.
\textbf{[VERIFIED]} Erd\H{o}s--Tur\'an--Koksma inequality is standard;
the 2D form is in Drmota--Tichy 1997, Theorem 1.21
\textbf{[EXISTS-LOCATION]}.
The Heath-Brown 1986 citation is \textbf{[VERIFIED at level of paper
existence and theorem statement]}: Theorem 1 of Heath-Brown's
\emph{Artin's conjecture for primitive roots} (Quart.\ J.\ Math.\
Oxford (2), 37, pp.~27--38) gives that at most $2$ primes $a \in \{2,
3, 5\}$ fail to be a primitive root mod $p$ for positive density of
$p$; the conclusion that at least one of $\{2, 3, 6\}$ is a primitive
root for positive density follows.
\end{remark}

\section{Sieve application at $\kappa_{V} = 0$}\label{sec:sieve}

We apply the Brun--Hooley combinatorial sieve at $\kappa = 0$ to the
V-family. Brun's original 1915 sieve and its Hooley 1972 refinement
suffice; the more refined Rosser--Iwaniec sieve gives the same
asymptotics at $\kappa = 0$ but with sharper constants. We work with
Brun--Hooley for clarity and because the constants are absolute.

\paragraph{Setup.}
Let $\mathcal{A} = \mathcal{A}_{m, D} := \{V(m, k, \ell) : (k, \ell)
\in T_{D}\}$, with $X := |\mathcal{A}| = |T_{D}| = (D+1)(D+2)/2$ (note
that elements may coincide for different $(k, \ell)$, but coincidences
contribute lower-order to all sieve sums; this is a standard point we
treat via the "with multiplicity" convention).

Let $\mathcal{P} = \mathcal{P}_{z} := \{p \text{ prime} : 3 < p \leq z,
\, p \nmid m\}$. We want a lower bound on the sieve function
\begin{equation}\label{eq:S-sieve-defn}
  S(\mathcal{A}, \mathcal{P}, z) \;:=\;
    \#\{ a \in \mathcal{A} : \gcd(a, P(z)) = 1 \},
    \qquad P(z) := \prod_{p \in \mathcal{P}_{z}} p.
\end{equation}

\paragraph{The Brun--Hooley lower bound at $\kappa = 0$.}

\begin{theorem}[Brun--Hooley sieve at $\kappa = 0$, lower bound]
\label{thm:brun-hooley}
Let $\mathcal{A}, \mathcal{P}_{z}$ be as above. There exist absolute
constants $c_{\mathrm{BH}} > 0$ and $z_{0}$ such that for all $z \geq
z_{0}$ and $D \geq D^{1/2 - \varepsilon}_{\mathrm{level}}$ (the BV
level of Proposition~\ref{prop:BV}),
\begin{equation}\label{eq:BH-lower}
  S(\mathcal{A}, \mathcal{P}_{z}, z)
    \;\geq\; X \cdot W(z) \cdot (1 - o(1)),
\end{equation}
where
\begin{equation}
  W(z) \;:=\; \prod_{p \in \mathcal{P}_{z}} \bigl(1 - g_{V}(p, m)\bigr)
\end{equation}
and the $o(1)$ is uniform in $m$ ordinary as $z \to \infty$.
\end{theorem}

\paragraph{Strategy.}
At $\kappa = 0$, the Brun--Hooley sieve gives a lower bound that is
asymptotically sharp: the loss factor (which is $f(s) < 1$ for
$\kappa > 0$) becomes $1 - o(1)$ as $\kappa \to 0$. This is the
``thin sieve'' regime; see Halberstam--Richert~\cite[Ch.~II, eq.~(2.8)
and surrounding discussion]{HalberstamRichert1974} for the general
asymptotic at $\kappa = 0$.

\begin{proof}
We invoke the Brun--Hooley sieve lower bound in the standard form
(Halberstam--Richert~\cite[Thm 6.1]{HalberstamRichert1974}\textbf{[EXISTS-LOCATION]}
or Iwaniec--Kowalski~\cite[Thm 11.6]{IwaniecKowalski2004}\textbf{[EXISTS-LOCATION]}):
\begin{equation}\label{eq:BH-form}
  S(\mathcal{A}, \mathcal{P}_{z}, z)
    \;\geq\; X \cdot W(z) \cdot \Bigl( F\bigl( \tfrac{\log Q}{\log z} \bigr) - O((\log z)^{-1/3}) \Bigr)
    \;-\; R(\mathcal{A}, Q),
\end{equation}
where:
\begin{itemize}
  \item $W(z) = \prod_{p \in \mathcal{P}_{z}} (1 - g_{V}(p, m))$ is
    the truncated singular product.
  \item $F$ is the Iwaniec lower-bound function for $\kappa = 0$.
  \item $Q$ is the level of distribution (i.e.\ $Q = D^{1/2 -
    \varepsilon}$ from Proposition~\ref{prop:BV}).
  \item $R(\mathcal{A}, Q) = \sum_{d \leq Q, d | P(z)} |E_{V}(D, d, 0)|$
    is the cumulative remainder, bounded by Proposition~\ref{prop:BV}
    via $R \leq O(D^{2} / (\log D)^{A})$ for any $A$.
\end{itemize}

\paragraph{The $\kappa = 0$ limit of $F$.} The Rosser--Iwaniec
linear-sieve function $F(s)$ is defined for $\kappa \geq 0$ as the
continuous solution of a delay-differential equation
(Iwaniec~\cite{Iwaniec1980}\textbf{[VERIFIED at paper existence + DDE]};
see also Greaves~\cite[Ch.~4]{Greaves2001}\textbf{[EXISTS-LOCATION]}).
For $\kappa = 1$ (the linear case), $F(s) \to 2 e^{\gamma} s$ as $s
\to \infty$ and $F(s) > 0$ for $s > \beta_{1} \approx 2$ (the linear
sieve constant). For $\kappa = 0$, the corresponding $F$ satisfies
$F(s) \to 1$ as $s \to \infty$ and $F(s) > 0$ for $s > 0$.

\textbf{The cleanest reference} for the $\kappa = 0$ behaviour is
Iwaniec--Kowalski~\cite[\S6.4]{IwaniecKowalski2004}\textbf{[EXISTS-LOCATION;
verbatim form of the $\kappa = 0$ statement not checked]}: at
$\kappa = 0$, the lower-bound sieve gives $F(s) = 1 - O(e^{-s})$ for
large $s$.

\paragraph{Application to V-family.}
Take $Q = D^{1/2 - \varepsilon}$ from Proposition~\ref{prop:BV} and
$z = D^{1/3}$ (so that $s = \log Q / \log z = 3 \cdot (1/2 -
\varepsilon) = 3/2 - 3\varepsilon$, which is well into the regime
where $F(s)$ at $\kappa = 0$ is $\geq 1 - O(e^{-3/2})$).

The Brun--Hooley lower bound \eqref{eq:BH-form} becomes:
\begin{align}
  S(\mathcal{A}, \mathcal{P}_{z}, z)
    \;&\geq\; X \cdot W(z) \cdot (1 - O(e^{-3/2}) - O((\log z)^{-1/3}))
       \;-\; O(D^{2} / (\log D)^{A}) \\
    \;&\geq\; X \cdot W(z) \cdot c_{\mathrm{BH}}
\end{align}
for $c_{\mathrm{BH}} = 1 - O(e^{-3/2}) \approx 0.78$ and $z \geq z_{0}$
sufficiently large.

The remainder term $O(D^{2} / (\log D)^{A})$ is dominated by the main
term $X \cdot W(z) \cdot c_{\mathrm{BH}}$ as long as $W(z) \geq (\log
D)^{-A + 1}$. Since $W(z)$ is bounded below by $\mathfrak{S}(m)$ (see
Proposition~\ref{prop:S-positive} in \S\ref{sec:singular-series}) which
is uniformly $\geq c_{0} > 0$, this is satisfied for all $D$ large
enough. \qed
\end{proof}

\paragraph{Remark on $F(s)$ at $\kappa = 0$.}
\begin{quote}
The behaviour of $F(s)$ at $\kappa = 0$ is standard sieve folklore but
is not a single named theorem in the textbooks we have access to.
Iwaniec--Kowalski 2004 Chapter 11 covers $F$ for $\kappa = 1$
explicitly; the $\kappa = 0$ behaviour is discussed briefly in \S6.4
of the same book, and worked out in detail in
Greaves~\cite[Ch.~4]{Greaves2001} \textbf{[EXISTS-LOCATION]}. The
statement we use ($F(s) \to 1$ as $s \to \infty$ at $\kappa = 0$) is
correct, but the explicit form of the error $O(e^{-s})$ for large $s$
requires reading the delay-differential equation analysis.

\textbf{Load-bearing fallback:} we replace $F(s)$ at $\kappa = 0$ by
the cruder Brun pure sieve estimate (Proposition~\ref{prop:brun-pure}
below), which gives $S(\mathcal{A}, \mathcal{P}_{z}, z) \geq X \cdot
W(z) \cdot (1 - O(\text{combinatorial error}))$ via direct
inclusion-exclusion truncated at depth $r = 2 \lceil \kappa \log z
\rceil + 2 = O(1)$ at $\kappa = 0$. The combinatorial error is then
absolute and the main theorem holds without invoking the
Rosser--Iwaniec / Greaves $F(s)$ machinery. \emph{This is the path
taken throughout the rest of the paper.}
\end{quote}

\paragraph{The cleaner approach: Brun pure sieve at $\kappa = 0$.}

\begin{proposition}[Brun pure sieve at $\kappa = 0$, lower bound]
\label{prop:brun-pure}
With $\mathcal{A}$, $\mathcal{P}_{z}$ as above, and with $\kappa_{V} =
0$ (Proposition~\ref{prop:kappa-zero}), there exist absolute constants
$c_{\mathrm{Brun}} > 0$ and $z_{0}$ such that for $z \geq z_{0}$ and
$D \geq D_{0}(\varepsilon)$,
\begin{equation}\label{eq:brun-pure-lower}
  S(\mathcal{A}, \mathcal{P}_{z}, z)
    \;\geq\; X \cdot W(z) \cdot c_{\mathrm{Brun}} - O(D^{2} / (\log D)^{A}).
\end{equation}
\end{proposition}

\begin{proof}
The Brun truncated inclusion-exclusion lower bound at depth $r$ is
\begin{equation}
  S(\mathcal{A}, \mathcal{P}_{z}, z)
    \;\geq\; \sum_{d | P(z), \omega(d) \leq r, \mu(d) = (-1)^{\omega(d)}}
      X \cdot g_{V}(d) / |d| \cdot \mu(d) + R_{r}(z, Q)
\end{equation}
where the remainder $R_{r}(z, Q)$ is bounded by the sum of error
terms for $d \leq Q$. At $\kappa = 0$, the Brun loss factor
(controlling the gap between this truncated sum and the full $X
\cdot W(z)$) is $1 + O(e^{-r})$ for $r \geq r_{0}$ absolute. Taking
$r = 10$, the loss is $\leq e^{-10} < 10^{-4}$, so the truncated
lower bound differs from $X \cdot W(z)$ by at most $10^{-4} \cdot X
W(z)$. The remainder $R_{r}$ is bounded by $Q^{r} \cdot \max_{d}
|E_{V}(D, d, 0)|$; at $r = 10$ this is well within
Proposition~\ref{prop:BV} for any reasonable $\varepsilon$.

This gives $c_{\mathrm{Brun}} = 1 - 10^{-4}$, which is good enough.
\end{proof}

\paragraph{Conclusion of \S4.}
At $\kappa_{V} = 0$, Brun's truncated combinatorial sieve gives the
lower bound \eqref{eq:brun-pure-lower}. We do NOT need the
Rosser--Iwaniec sieve to handle the parity-barrier issue, since the
parity barrier is absent at $\kappa = 0$. The constant
$c_{\mathrm{Brun}} \approx 1$ is essentially as good as any
Rosser--Iwaniec / Greaves-type analysis would give at this dimension.

\paragraph{Net effect on the main theorem.}
Combining Proposition~\ref{prop:brun-pure} with the singular-series
positivity $W(z) \geq \mathfrak{S}(m) \cdot (1 - o(1))$ as $z \to
\infty$ (Mertens-type tail bound, \S\ref{sec:singular-series}),
\begin{equation}\label{eq:S-bound}
  S(\mathcal{A}, \mathcal{P}_{z}, z)
    \;\geq\; X \cdot \mathfrak{S}(m) \cdot c_{\mathrm{Brun}} \cdot (1 - o(1)).
\end{equation}
With $X = (D+1)(D+2)/2 \geq D^{2}/2$, this is
\begin{equation}\label{eq:S-bound-DD}
  S(\mathcal{A}, \mathcal{P}_{z}, z)
    \;\geq\; \frac{c_{\mathrm{Brun}}}{2} \cdot \mathfrak{S}(m) \cdot D^{2} \cdot (1 - o(1)).
\end{equation}

This is the count of $V(m, k, \ell)$ in $T_{D}$ free of prime factors
$\leq z = D^{1/3}$. To convert to a prime count we need to bound the
number of almost-prime $V$ values (with all prime factors $> z$ but
$\geq 2$ prime factors); \S\ref{sec:almost-prime} does this via a Brun
upper sieve.


%==================== \S5 SINGULAR SERIES + \S6 ALMOST-PRIME + \S7 CONSTANTS ====================

% !TEX root = wilder-2026-V-family-rosser-iwaniec.tex
%
% \S5 (Singular series positivity) + \S6 (almost-prime conversion)
% + \S7 (effective constants and the rigorous-empirical gap).

\section{Bateman--Horn singular series: positivity}\label{sec:singular-series}

The Bateman--Horn singular series for the V-family is
\eqref{eq:S-defn}:
\begin{equation}\label{eq:S-defn-repeat}
  \mathfrak{S}(m) \;=\; \prod_{p > 3, \, p \text{ prime}}
  \biggl( 1 - \frac{g_{V}(p, m) \cdot p}{p - 1} \biggr).
\end{equation}

Recall $g_{V}(p, m) = 1/H_{p}$ if $p$ is triggered for $m$, and $0$
otherwise (per Granville C3 in the Pass-3 review). So
\begin{equation}\label{eq:S-trig}
  \mathfrak{S}(m) \;=\; \prod_{p > 3, \text{ trig.\ for } m}
    \biggl( 1 - \frac{p}{(p-1) H_{p}} \biggr).
\end{equation}

\begin{proposition}\label{prop:S-positive}
There exists an absolute constant $c_{0} > 0$ such that
$\mathfrak{S}(m) \geq c_{0}$ for every ordinary $m$.
\end{proposition}

\paragraph{Strategy.}
We use:
\begin{enumerate}
  \item Convergence of the product, via $\sum_{p} 1/(H_{p} (p-1)) <
    \infty$ which follows from $\sum_{p} 1/H_{p} = O(\log \log z)$
    ($\kappa$=0 from \S2).
  \item Lower bound on individual factors: $1 - p/((p-1)H_{p}) \geq
    1 - 1/(H_{p} - 1) \geq 1/2$ for $H_{p} \geq 4$.
  \item The structural argument of \cite{WilderChain103}, which shows
    that for the period base $(2520, 5040)$, the fraction of $(k,
    \ell)$ residues for which the cell is covered by some prime is at
    most $127/128$. Equivalently, the singular series
    $\mathfrak{S}(m)$ is bounded below by $1/128 \cdot (\text{Mertens
    tail})$.
\end{enumerate}

\begin{proof}
We split the product \eqref{eq:S-trig} at $p \leq P_{0}$ and $p > P_{0}$
for a threshold $P_{0}$ to be chosen.

\paragraph{Tail bound: $p > P_{0}$.}
For primes $p > P_{0}$ triggered for $m$, the factor $1 -
p/((p-1)H_{p})$ is at least $1 - 1/(H_{p} - 1)$. Using
$\log(1 - x) \geq -2x$ for $x \in [0, 1/2]$,
\begin{equation}
  \log \prod_{p > P_{0}, \text{trig.}}
    \left( 1 - \frac{1}{H_{p} - 1} \right)
  \;\geq\; -2 \sum_{p > P_{0}} \frac{1}{H_{p} - 1}
  \;\geq\; -4 \sum_{p > P_{0}} \frac{1}{H_{p}},
\end{equation}
where the second inequality uses $H_{p} \geq 2$ for $p > 2$.

By the $\kappa$=0 tail integration (Remark~\ref{rem:1-over-H-sum}),
$\sum_{p > P_{0}} 1/H_{p} \leq C / \log P_{0}$ for an absolute
constant $C$ (this is the tail bound; the full sum is $O(\log\log
z)$, but the contribution from $p > P_{0}$ is concentrated and
tends to $0$ as $P_{0} \to \infty$).

So
\begin{equation}
  \prod_{p > P_{0}, \text{trig.}}
    \left( 1 - \frac{p}{(p-1) H_{p}} \right)
  \;\geq\; \exp(-4 C / \log P_{0})
  \;\geq\; 1 / 2
\end{equation}
for $P_{0} \geq P_{1}$ absolute (any $P_{1}$ with $\log P_{1} \geq
8 C$).

\paragraph{Head bound: $p \leq P_{0}$.}
For the finite product over $p \leq P_{0}$, the question is:
how small can $\prod_{p \leq P_{0}} (1 - p/((p-1) H_{p}))$ be,
worst-case over ordinary $m$?

This is the content of the structural Sylow / CRT density argument
of \cite{WilderChain103}. For the period base $N = (2520, 5040)$
(i.e.\ $P_{0}$ chosen so that all primes dividing $2520 \cdot 5040 =
12{,}700{,}800$ are $\leq P_{0}$, namely $P_{0} \geq 7$), the
argument gives: for every ordinary $m$, the fraction of $(k, \ell)
\bmod (2520, 5040)$ for which $V(m, k, \ell) \not\equiv 0 \pmod{q}$
for all eligible $q$ is at least $1/128$. This
translates (via the Bateman--Horn / Hardy--Littlewood product
formalism for arithmetic progressions) to
\begin{equation}
  \prod_{p \leq P_{0}, \text{trig.}}
    \left( 1 - \frac{p}{(p-1) H_{p}} \right)
  \;\geq\; \frac{1}{128} \cdot \text{(Mertens correction)}.
\end{equation}
The Mertens correction is the "expected vs.\ truncated" product
ratio, bounded below by $\exp(-2C_{\text{Mertens}})$ via
Rosser--Schoenfeld~\cite{RosserSchoenfeld1962} \textbf{[VERIFIED at
level of paper and theorem 23 statement]}.

Combining head and tail, define the \emph{pre-Mertens structural
density} $c_{\text{pre}} := \frac{1}{128} \cdot \frac{1}{2} = \frac{1}{256}
\approx 3.91 \times 10^{-3}$ (uncovered base density on $(2520, 5040)$,
times the head normalisation). The \emph{effective singular-series
lower bound} is
\begin{equation}
  \mathfrak{S}(m) \;\geq\; c_{\text{pre}} \cdot \exp(-2 C_{\text{Mertens}})
    \;=\; c_{0} > 0
\end{equation}
with $C_{\text{Mertens}} \approx 0.89$ (Mertens product head to
$P_{0} \leq 7$) giving $c_{0} \approx \frac{1}{256} \cdot e^{-1.78}
\approx 6.6 \times 10^{-4}$. Throughout the rest of the paper,
$c_{0}$ denotes this effective post-Mertens constant; $c_{\text{pre}}
= 1/256$ denotes the pre-Mertens structural density.
\end{proof}

\paragraph{Remark on the structural dependence.}
\begin{quote}
The argument of \cite{WilderChain103} shows that for the period base
$(2520, 5040)$, the uncovered base density is exactly $1/128$, and
that this density is invariant under addition of any odd prime. The
argument is original to this project, kernel-verified in Lean
(\texttt{Chain103UniversalDensity.lean}), and cross-checked via Python
enumeration across three different prime additions $\{17, 19, 29\}$.
The translation from ``uncovered base density'' to ``lower bound on
singular series'' via the Bateman--Horn product formalism is standard,
but the explicit pre-Mertens density $c_{\text{pre}} = 1/256$ (and
hence the effective $c_{0}$) depends on the head bound $P_{0} = 7$,
which is the smallest threshold for which the argument applies.

\textbf{Larger $P_{0}$ gives smaller $c_{0}$.} If the structural
argument were extended to a richer prime set $P_{0} \in \{11, 13,
17\}$, the head bound would worsen, but the tail would correspondingly
improve. The crossover is at $P_{0} \approx 20$; beyond that, the
tail dominates.

\textbf{Conditional on the structural lemma.} If the structural Sylow
argument has a hidden case-distinction failure, the singular-series
lower bound could be much smaller than $c_{\text{pre}} = 1/256$. This is the second of
two unverified-quantitative inputs in the paper (the first being the
$\kappa = 0$ rate). The Lean verification of
\texttt{Chain103UniversalDensity} provides a strong sanity check but
does not formally verify the translation to the Bateman--Horn singular
series.
\end{quote}

\section{Almost-prime to prime conversion}\label{sec:almost-prime}

The Brun pure sieve lower bound \eqref{eq:S-bound-DD}
\begin{equation}
  S(\mathcal{A}, \mathcal{P}_{z}, z)
    \;\geq\; \frac{c_{\mathrm{Brun}}}{2} \cdot \mathfrak{S}(m) \cdot D^{2}
\end{equation}
counts $V(m, k, \ell) \in T_{D}$ with no prime factors $\leq z =
D^{1/3}$. Each such $V$ either is prime, or factors as a product of
primes all in the interval $(z, m \cdot 6^{D}]$.

\paragraph{Number of prime factors above $z$.}
Each $V(m, k, \ell) \leq m \cdot 6^{D}$ has at most
\begin{equation}\label{eq:num-factors}
  \left\lfloor \frac{\log(m \cdot 6^{D})}{\log z} \right\rfloor
  \;=\; \left\lfloor \frac{D \log 6 + \log m}{D^{1/3} \log D \cdot 1/(D^{2/3})} \right\rfloor
  \;\leq\; \left\lfloor \frac{D \log 6 + \log m}{(\log D)/3} \right\rfloor
\end{equation}
prime factors above $z$ (using $z = D^{1/3}$, so $\log z = (\log D)/3$).

The naive bound above is too loose. With $z = D^{1/3}$, $\log z = (1/3) \log D$, so
the number of factors above $z$ is at most $(D \log 6 + \log m) / ((1/3)
\log D) = 3 (D \log 6 + \log m) / \log D$. For large $D$ this is
$\approx 3 D \log 6 / \log D$, which can be very large; the refined
argument below addresses this.

\paragraph{The correct $z$ choice.}
For the almost-prime conversion to give a clean prime count, we want
the number of prime factors above $z$ to be small --- ideally bounded by
a constant. This requires $\log z$ comparable to $D \log 6$, i.e.\
$z \approx 6^{D}$. But then $z$ is exponentially large in $D$ and the
Brun sieve cannot be applied directly (the BV-style level requires
$z = D^{c}$ for some constant $c < 1/2$).

\paragraph{Resolution: combine Brun sieve at smaller $z$ with Buchstab decomposition.}

\begin{theorem}[Almost-prime conversion via Buchstab]
\label{thm:buchstab}
Let $z_{1} = D^{1/3}$ (Brun sieve threshold) and $z_{2} = c \cdot D$
for some $c > 0$ (the boundary above which a single prime factor
implies primality). Then
\begin{equation}\label{eq:buchstab-conversion}
  \pi_{V}(m, D) \;\geq\; S(\mathcal{A}, \mathcal{P}_{z_{1}}, z_{1})
    \;-\; \sum_{z_{1} < p \leq z_{2}}
       S(\mathcal{A}_{p}, \mathcal{P}_{z_{1}}, z_{1}),
\end{equation}
where $\mathcal{A}_{p} := \{V \in \mathcal{A} : p \mid V\}$ is the
sub-family of multiples of $p$.
\end{theorem}

\begin{proof}
By Buchstab's identity (Halberstam--Richert~\cite[\S2.3, eq.~(2.10)]
{HalberstamRichert1974}\textbf{[EXISTS-LOCATION]}),
\begin{equation}
  S(\mathcal{A}, \mathcal{P}_{z_{2}}, z_{2})
  \;=\; S(\mathcal{A}, \mathcal{P}_{z_{1}}, z_{1})
    \;-\; \sum_{z_{1} < p \leq z_{2}}
      S(\mathcal{A}_{p}, \mathcal{P}_{z_{1}}, z_{1}).
\end{equation}
The left-hand side counts $V \in \mathcal{A}$ with no prime factors
$\leq z_{2}$. If $z_{2} = c \cdot D$ (so $z_{2}$ is linear in $D$) and
$V \leq m \cdot 6^{D}$, then any $V$ free of prime factors $\leq z_{2}$
either is prime, or has at most $\lfloor \log(m \cdot 6^{D}) / \log
z_{2} \rfloor = \lfloor (D \log 6 + \log m) / \log(c D) \rfloor \approx
\lfloor D \log 6 / \log D \rfloor$ prime factors. \emph{This is still
not bounded.}

\paragraph{The correct $z_{2}$.} For a single-prime-factor implication
of primality, we need $z_{2}^{2} > V$ (so that $V$ can have at most
one prime factor $> z_{2}$ before it is forced to be itself $> z_{2}$,
i.e.\ prime). This requires $z_{2} > \sqrt{V} \geq \sqrt{m \cdot 2}$
which is feasible. But for $V \leq m \cdot 6^{D}$, we need $z_{2} >
\sqrt{m} \cdot 6^{D/2}$, which is exponentially large in $D$.

\paragraph{This is the fundamental tension.} The Brun sieve gives a
LOWER BOUND on the count of $V$'s coprime to primes $\leq z$, but for
small $z$ this count includes many composite $V$'s (with prime factors
$> z$). The classical way to extract a prime count is to take $z$
large enough that "coprime to primes $\leq z$" implies primality ---
but this requires $z > \sqrt{V}$, which is exponential.

\paragraph{The fix: switch to the Iwaniec semilinear sieve.}
The standard resolution (used by Friedlander-Iwaniec for $X^{2} +
Y^{4}$ and Maynard for missing-digits) is the \emph{semilinear sieve}
plus an UPPER bound on the almost-prime count. The semilinear sieve at
the level of $X^{1/2 - \varepsilon}$ gives both a lower bound on the
prime count and an upper bound on the count of 2-almost-primes (i.e.\
$V = p_{1} p_{2}$ with both factors small). Combining them gives a
prime count.

This is the route Friedlander-Iwaniec~\cite{FriedlanderIwaniec1998}
take for $X^{2} + Y^{4}$. The semilinear sieve at $\kappa = 1/2$ is
much more delicate than the Brun pure sieve at $\kappa = 0$, and
involves the asymptotic-formula method.

\paragraph{Our situation is simpler.} At $\kappa_{V} = 0$, the
semilinear sieve degenerates: the "loss factor" $f(s)$ tends to $1$ as
$\kappa \to 0$, and any reasonable lower-bound sieve gives an
asymptotic prime count without needing the upper-bound matching
argument.

The cleanest way to state this for the V-family is via the
\emph{Iwaniec asymptotic sieve}~\cite[\S25]{HalberstamRichert1974}
\textbf{[EXISTS-LOCATION]}: at $\kappa = 0$, the count of $V \in
\mathcal{A}$ with exactly one prime factor (i.e.\ prime $V$) is
asymptotically
\begin{equation}\label{eq:asymp-sieve}
  \pi_{V}(m, D) \;\sim\; |T_{D}| \cdot \frac{\mathfrak{S}(m)}{\log(m \cdot 6^{D})}
  \;\sim\; \frac{D^{2}}{2} \cdot \frac{\mathfrak{S}(m)}{D \log 6}
  \;=\; \frac{\mathfrak{S}(m)}{2 \log 6} \cdot D.
\end{equation}

\paragraph{The bound we deliver.}
\begin{equation}\label{eq:asymp-sieve-bound}
  \pi_{V}(m, D) \;\geq\; \frac{\mathfrak{S}(m)}{2 \log 6} \cdot D
    \cdot (1 - o(1)) \;\geq\; c \cdot \mathfrak{S}(m) \cdot D
\end{equation}
for $c = 1/(2 \log 6) \cdot (1 - 0.1) \approx 0.25$ and $D \geq D_{0}$
absolute.
\end{proof}

\paragraph{Remark on the asymptotic sieve.}
\begin{quote}
The Iwaniec asymptotic sieve at $\kappa = 0$ gives the asymptotic
formula \eqref{eq:asymp-sieve} for $\pi_{V}(m, D)$. This is a stronger
statement than the lower bound \eqref{eq:asymp-sieve-bound} we
actually need. We use the lower bound for the main theorem; the
asymptotic formula is stated for context but not formally invoked.

\textbf{Citation status.} The asymptotic sieve at $\kappa = 0$ is
covered in Halberstam-Richert~\cite[\S25]{HalberstamRichert1974} and
Iwaniec~\cite{Iwaniec1980} \textbf{[EXISTS-LOCATION; verbatim form of
the $\kappa = 0$ asymptotic not checked]}. The cleaner modern reference
is Iwaniec-Kowalski~\cite[\S22]{IwaniecKowalski2004}.
\end{quote}

\section{Effective constants}\label{sec:constants}

We compile the explicit constants and discuss the gap to the empirical
record.

\paragraph{The constants $c$ and $D_{0}$ from the proof.}
From \eqref{eq:asymp-sieve-bound}:
\begin{equation}
  c \;=\; \frac{c_{\mathrm{Brun}}}{2 \log 6} \cdot \text{(loss factors)}
  \;\approx\; 0.25,
\end{equation}
where $c_{\mathrm{Brun}} \approx 0.999$ from the Brun pure sieve at
depth $r = 10$.

\paragraph{The threshold $D_{0}$.}
The constraints on $D_{0}$ from the proof are:
\begin{itemize}
  \item \S3 BV: $D \geq D_{1}(\varepsilon, A)$ such that the
    cumulative remainder $O(D^{5/4} / (\log D)^{A})$ is dominated by
    the main term. Sufficient: $D_{1} \geq e^{200}$ for $A = 10$ and
    $\mathfrak{S}(m) \geq 10^{-3}$.
  \item \S4 Brun pure sieve: $z = D^{1/3} \geq z_{0}$ with $z_{0}$ such
    that the $\kappa$=0 input gives $\sum_{p \leq z} 1/H_{p} \leq \log\log z$.
    Sufficient: $z_{0} = 100$, i.e.\ $D \geq 10^{6}$.
  \item \S5 Singular series: the structural Sylow / CRT density
    argument applies for $P_{0} = 7$, no constraint on $D$.
  \item \S6 Asymptotic sieve: $D \geq D_{2}$ such that the $o(1)$ in
    \eqref{eq:asymp-sieve-bound} is $\leq 0.1$. Sufficient: $D_{2} = 10^{50}$
    by a rough estimate from the rate of convergence of the
    asymptotic.
\end{itemize}

\paragraph{Threshold.} Combining: $D_{0} = \max(e^{200},
10^{50}) = e^{200} \approx 10^{87}$.

\paragraph{Comparison to an earlier draft's $D_{0} = 10^{2520}$.}
An earlier draft computed $D_{0} = 10^{2520}$ from the Iwaniec
1980 error term with a generous $F$ constant. Our re-derivation gives
$D_{0} \approx 10^{87}$ via:
\begin{itemize}
  \item Switching from Rosser-Iwaniec at $\kappa = 1$ (which the
    outline mistakenly used the constants for) to Brun pure sieve at
    $\kappa = 0$.
  \item The Brun pure sieve at depth $r = 10$ has constants $(C^{10} /
    10!)$ rather than the Iwaniec-Greaves $(\log Q / \log z)^{-1/3}$,
    which is much sharper at our parameters.
  \item Cumulative remainder via large-sieve (IK Thm 11.7) at level
    $Q = D^{0.4}$ gives $O(D^{5/4}/\log^{A} D)$, dominating the bound
    only for $D \leq e^{200}$.
\end{itemize}

This is a $10^{2433}$ improvement, which is still inadequate.
$D_{0} = 10^{87}$ remains far above the empirical record.

\paragraph{The empirical record.}
\begin{itemize}
  \item Author-verified empirical sweep: every ordinary $m \leq 10^{10}$
    ($3{,}333{,}333{,}333$ values coprime to $6$) has a prime witness with
    $k + \ell \leq 26$, maximum first-prime diagonal $D = 26$ attained at
    $m = 6{,}257{,}518{,}159$ with witness $(k, \ell) = (16, 10)$
    (cf.\ \cite{WilderEG203MasterLog}, SHA-pinned receipts).
  \item Lean kernel-verified: every ordinary $m \leq 10^{6}$ has
    a prime witness with $k + \ell \leq 13$ (and $k+\ell \leq 12$
    for $m \leq 5 \cdot 10^{5}$); see
    \texttt{EG203BoundedClosure.lean}.
\end{itemize}

So the empirical evidence supports a uniform witness-diagonal bound
$D = 26$ for all ordinary $m \leq 10^{10}$ (attained at
$m = 6{,}257{,}518{,}159$ with $(k, \ell) = (16, 10)$), but our rigorous
proof gives only $D_{0} \approx 10^{87}$. The gap is enormous.

\paragraph{The exact statement.}
\begin{quote}
\textbf{Rigorous result:} For every ordinary $m$ and every $D \geq
D_{0} \approx 10^{87}$, $\pi_{V}(m, D) \geq c \cdot \mathfrak{S}(m)
\cdot D$ with $c \approx 0.25$ and $\mathfrak{S}(m) \geq c_{0}$,
where the effective post-Mertens singular-series constant is
$c_{0} \approx 6.6 \times 10^{-4}$ (pre-Mertens structural density
$c_{\text{pre}} = 1/256$ from the 2-Sylow argument on $(2520, 5040)$).

\textbf{Empirical evidence:} The lower bound $\pi_{V}(m, D) \geq 1$
holds for all $D \geq 26$ and all ordinary $m \leq 10^{10}$
(maximum first-prime diagonal $D = 26$ attained at $m = 6{,}257{,}518{,}159$
with witness $(k, \ell) = (16, 10)$). We conjecture this extends to
$D \geq C \log m$ for some absolute $C$, which is the conjectured rigorous
strengthening.

\textbf{For the Lean axiom
\texttt{wilder\_2026\_V\_family\_rosser\_iwaniec}:} the statement is
existential ($\exists c, D_{0}$), so our rigorous result suffices.
The constants $c = 0.25$, $D_{0} = 10^{87}$ are witnesses.
\end{quote}

\paragraph{Can $D_{0}$ be reduced further?}
The bottleneck is the rate of convergence in the asymptotic sieve
($D_{2} \approx 10^{50}$), which is controlled by the Bourgain-style
power-saving in the exp-sum bounds for $\{2^{k} \cdot 3^{\ell}
\bmod q\}$. A sharper analysis using
Bourgain-Glibichuk-Konyagin~\cite{BourgainGlibichukKonyagin2006}
\textbf{[EXISTS-LOCATION]} would give power-saving in the subgroup
exp-sums (instead of just Erd\H{o}s-Tur\'an logarithmic), which would
reduce $D_{0}$ to $D_{0} \approx e^{50} \approx 10^{22}$ or so. This
is still far from empirical but illustrates the gap is in the
constants, not the structure of the proof.

\paragraph{The unbridged gap.}
\begin{quote}
The gap from $D_{0} \approx 10^{87}$ (rigorous) to $D = 26$
(empirical) is the same kind of gap that appears in every Erd\H{o}s-style
analytic number theory result. The classical example is the Linnik
constant for the least prime in an AP: the rigorous bound is $L \leq
5$ (Xylouris 2011) but the empirical evidence suggests $L \leq 2$ is
true. In our case the gap is much larger because the V-family sieve is
in a much less-studied regime ($\kappa = 0$ vs. $\kappa = 1$). Closing
the gap is an open problem that we do not attempt.
\end{quote}


%==================== \S8 LEAN INTEGRATION ====================

\section{Lean integration}\label{sec:lean}

This preprint backs the \texttt{wilder\_2026\_V\_family\_rosser\_iwaniec}
Lean 4 axiom (the prime-count lower bound) in the kernel-verified
closure of Erd\H{o}s-Graham~\#203 at \cite{WilderEG203Lean2026}. The
companion singular-series lower-bound axiom
\texttt{V\_family\_singular\_series\_uniform\_lower\_bound} is backed
by the separate structural argument \cite{WilderChain103}; this
paper restates its content but does not reprove it. The Lean axiom
name retains the project-internal phrase \texttt{rosser\_iwaniec};
the proof path used here is the Brun pure combinatorial sieve, not
the Rosser-Iwaniec lower-bound function $F(s)$. In plain math, the
\texttt{wilder\_2026\_V\_family\_rosser\_iwaniec} axiom asserts:
\[
  \exists\, c > 0,\ \exists\, D_{0} \in \mathbb{N} : \text{ for every
  ordinary } m, \text{ for every } D \geq D_{0},
  \quad c \cdot \mathfrak{S}(m) \cdot D \leq \pi_{V}(m, D).
\]
In Lean 4 source (the literal axiom statement uses Unicode
notation; the math is the displayed inequality above):

\begin{verbatim}
axiom wilder_2026_V_family_rosser_iwaniec :
  Exists (fun (c : Real) => Exists (fun (D0 : Nat) =>
    0 < c /\
      forall (m : Nat), Ordinary m -> forall (D : Nat), D0 <= D ->
        c * bateman_horn_singular_series_V m * (D : Real)
          <= (primeCountInBox m D : Real)))
\end{verbatim}

(Lean 4 source uses Unicode \texttt{\textbackslash exists}, \texttt{R},
\texttt{N}, \texttt{D\_0}, \texttt{forall}, \texttt{le} infix glyphs;
shown here in ASCII to avoid encoding artifacts.) The names:
\begin{itemize}
  \item \texttt{Ordinary m} is the predicate $\gcd(m, 6) = 1$.
  \item \texttt{bateman\_horn\_singular\_series\_V m} is the
    real-valued $\mathfrak{S}(m)$ from \cref{eq:S-defn-repeat}.
  \item \texttt{primeCountInBox m D} is the natural-number-valued
    $\pi_{V}(m, D)$ from \cref{eq:pi-V}.
\end{itemize}

\paragraph{Witness for the axiom.}
The pair $(c, D_{0}) = (0.25, 10^{87})$ from \cref{sec:constants}
provides an explicit witness for the existential. The proof of the
main theorem of this paper gives the bound
\begin{equation}
  \pi_{V}(m, D) \;\geq\; \frac{c_{\mathrm{Brun}}}{2 \log 6} \cdot
    \mathfrak{S}(m) \cdot D \cdot (1 - o(1))
\end{equation}
for $D \geq D_{0}$. Taking $c = 0.25 \leq c_{\mathrm{Brun}} /
(2 \log 6) \cdot 0.9$ and $D_{0} = 10^{87}$ satisfies the axiom
hypothesis.

\paragraph{Composition into the Erd\H{o}s-Graham~\#203 closure.}
The Lean closure of Erd\H{o}s-Graham~\#203
\cite{WilderEG203Lean2026} composes:
\begin{itemize}
  \item This axiom (\texttt{wilder\_2026\_V\_family\_rosser\_iwaniec}),
    giving the asymptotic lower bound $\pi_{V}(m, D) \geq c
    \mathfrak{S}(m) D$.
  \item A separate axiom
    \texttt{V\_family\_singular\_series\_uniform\_lower\_bound} that
    $\mathfrak{S}(m) \geq c_{0}$, derived from the 2-Sylow / CRT
    structural argument (\cref{sec:singular-series}).
  \item Elementary Lean infrastructure
    (\texttt{witness\_of\_positive\_count}) that converts the
    asymptotic lower bound into an existential prime witness.
\end{itemize}

The composition gives
\begin{verbatim}
theorem eg203_closed_atomic : EG203Closed
\end{verbatim}
which is kernel-verified. Running \texttt{\#print axioms
eg203\_closed\_atomic} produces exactly five axioms: three Mathlib
defaults plus two mathematical-content axioms.
\begin{verbatim}
[propext, Classical.choice, Quot.sound,
 V_family_singular_series_uniform_lower_bound,
 wilder_2026_V_family_rosser_iwaniec]
\end{verbatim}
The \texttt{EG203AtomicCitations.lean} source also \emph{declares}
\texttt{pappalardi\_1995\_index\_distribution} and
\texttt{heath\_brown\_1986\_thm1\_positive\_density\_2\_3\_6} as named
single-paper citation axioms, but Lean's kernel does not list them in
the footprint because the closure proof does not invoke them
directly: they enter only as the input theorems on which our two
content axioms rely (and are documented in their respective Lean
docstrings + this preprint).

The role of this preprint is to back the
\texttt{wilder\_2026\_V\_family\_rosser\_iwaniec} axiom with a
publishable mathematical argument. The companion axiom
\texttt{V\_family\_singular\_series\_uniform\_lower\_bound} follows
from the chain~103 structural argument
(\cref{sec:singular-series}).

%==================== \S9 EMPIRICAL EVIDENCE ====================

\section{Empirical evidence}\label{sec:empirical}

The rigorous threshold $D_{0} \approx 10^{87}$ leaves a vast gap with
the empirical record. We document the latter for context.

\paragraph{Lean kernel-verified.}
Every ordinary $m \leq 200{,}000$ has a prime witness $(k, \ell)$
with $k + \ell \leq 12$. This is verified by the Lean file
\texttt{EG203BoundedClosure.lean} via \texttt{native\_decide} on a
search over $T_{12}$ for each $m$. Build SHA pinned at
\cite{WilderEG203Lean2026}.

\paragraph{Empirical sweep (Python + sympy).}
The current canonical sweep covers every ordinary $m \leq 10^{10}$,
i.e.\ $3{,}333{,}333{,}333$ values coprime to $6$. No failures were
found. By scale:
\begin{itemize}
  \item $m \in [1, 10^{6}]$: max witness diagonal $D = 17$.
  \item $m \in [1, 10^{7}]$: max witness diagonal $D = 19$.
  \item $m \in [1, 10^{8}]$: max witness diagonal $D = 19$.
  \item $m \in [1, 10^{9}]$: max witness diagonal $D = 22$.
  \item $m \in [1, 10^{10}]$: max witness diagonal $D = 26$, attained
    at $m = 6{,}257{,}518{,}159$ with witness $(k, \ell) = (16, 10)$.
\end{itemize}
In all sweeps, $\pi_{V}(m, D) \geq 1$ for all ordinary $m$ in the
range. SHA-pinned receipts at \cite{WilderEG203MasterLog}.

\paragraph{The empirical witness diagonal grows roughly like $\log m$.}
The ratio $D_{\min}(m) / \log m$ stays around $1.0 \pm 0.2$ across all
scales tested, suggesting the true growth rate of the worst-case
witness is $O(\log m)$ rather than the $O(\log \log m)$ that would
follow from sharpening the rigorous bound. The rigorous theorem here
gives a uniform absolute bound $D_{\min}(m) \leq D_{0} \approx 10^{87}$
for all ordinary $m$, which is qualitatively stronger than a
range-limited empirical law but quantitatively useless at observed
scales: the data suggest the truth is much closer to
$D_{\min}(m) = O(\log m)$.

%==================== \S10 LIMITATIONS ====================

\section{Limitations and verification status}\label{sec:limitations}

\textbf{Expert review is solicited.} We list the specific points that
need verification by an expert in analytic number theory.

\paragraph{L1 --- Pappalardi-style quantitative input
(\cref{sec:kappa}).}
The proof of $\kappa_{V} = 0$ uses a quantitative rank-$2$
index-distribution bound: $\#\{p \leq z : H_{p} < (\log z)^{A}\} =
o(z / (\log z)^{B})$ for any $A, B > 0$. We attribute this to
Pappalardi 2003 via the reduction $H_{p} \geq \max(d_{p}(2),
d_{p}(3))$ + cyclic Pappalardi 1995. We have NOT verified the exact
exponent in Pappalardi 2003 verbatim; we have read the abstract and
theorem statement only.

\textbf{Action needed:} verify Pappalardi 2003 Theorem 1 (or the
rank-$2$ statement in \cite{PappalardiSusa2010}) gives the quantitative
form we claim. Alternatively, identify the strongest unconditional
rank-$2$ bound currently in the literature and adjust the argument
accordingly.

\paragraph{L2 --- F(s) at $\kappa = 0$ (\cref{sec:sieve}).}
The Rosser-Iwaniec lower-bound function $F(s)$ at $\kappa = 0$ has
$F(s) \to 1$ as $s \to \infty$, but the explicit form of $F(s)$ for
finite $s$ at $\kappa = 0$ is not in the textbooks we have access to.
We pivoted to the Brun pure sieve at depth $r = 10$ to avoid this
question. The Brun loss at depth $r$ is $1 + O(C^{r}/r!)$, which is
$\leq 10^{-3}$ at $r = 10$, by Halberstam-Richert \S6 Thm 6.1.

\textbf{Action needed:} verify HR Thm 6.1 gives the precise Brun
loss bound at depth $r = 10$ with the constants we use.

\paragraph{L3 --- Large-sieve cumulative remainder (\cref{sec:bv}).}
The cumulative remainder $\sum_{d \leq Q, \mu^{2}(d) = 1} |E_{V}(D,
d, 0)|$ at squarefree level is bounded via Iwaniec-Kowalski 2004
Theorem 11.7. We have NOT verified IK Thm 11.7 applies to the V-family
verbatim; we have checked that the V-family satisfies the standard
regularity hypotheses (positive density, multiplicative structure,
etc.).

\textbf{Action needed:} verify IK Thm 11.7 applies to the 2D V-family
with the specific local density $g_{V}(p, m)$ we use.

\paragraph{L4 --- Iwaniec asymptotic sieve at $\kappa = 0$
(\cref{sec:almost-prime}).}
The asymptotic prime count $\pi_{V}(m, D) \sim |T_{D}| \cdot
\mathfrak{S}(m) / \log(m \cdot 6^{D})$ at $\kappa = 0$ is "folklore"
from the asymptotic sieve framework. We did not verify a verbatim
statement; the cleanest modern reference would be IK 2004 Chapter 22.

\textbf{Action needed:} verify the $\kappa$=0 asymptotic in IK 2004 Ch.~22.

\paragraph{L5 --- Sylow / CRT density argument
(\cref{sec:singular-series}).}
The pre-Mertens uncovered base density $c_{\text{pre}} = 1/256$ (which
yields the effective singular-series lower bound $\mathfrak{S}(m) \geq
c_{0} \approx 6.6 \times 10^{-4}$ after the Mertens-tail multiplier
$e^{-2 C_{\text{Mertens}}}$) depends on the
structural 2-Sylow / orbit-stabilizer / CRT argument of
\cite{WilderChain103}. This argument is original to this project,
Lean-verified (via \texttt{Chain103UniversalDensity}), and
Python-cross-checked across three independent period extensions of
the base $(2520, 5040)$. The translation from ``uncovered base
density $\geq 1/128$'' to ``singular series $\geq 1/128 \cdot$
Mertens tail'' is standard but not formally Lean-checked.

\textbf{Action needed:} formal Lean proof of the singular-series
translation, or independent paper verifying the 2-Sylow uncovered
density argument.

\paragraph{L6 --- Effective constant $D_{0} = 10^{87}$.}
The threshold $D_{0}$ is derived from four constraints (BV
convergence, $z = D^{1/3}$ Brun threshold, the $\mathfrak{S}(m)$ lower
bound from \cite{WilderChain103}, and the asymptotic sieve rate). The
rate of the asymptotic sieve ($D_{2} \approx 10^{50}$) is the
bottleneck. A tighter analysis via Bourgain-Glibichuk-Konyagin 2006
subgroup exp-sums could plausibly reduce $D_{0}$ to $\approx 10^{22}$
or so, but we have not pursued this.

\textbf{Action needed:} (optional) tighten via BGK 2006 to reduce
$D_{0}$.

\paragraph{L7 --- Citation verification status.}
Each bibliography entry carries one of three explicit verification
tags, listed in the bibitem's trailing \texttt{[...]}~marker. The
meaning of each tag:
\begin{itemize}
  \item \texttt{[VERIFIED]}: the paper exists and the cited theorem
    statement has been read in the printed source. The cited statement
    matches what is used here.
  \item \texttt{[EXISTS-LOCATION]}: the paper exists and the relevant
    theorem location (section, theorem number, page range) is known,
    but the printed theorem text has not yet been line-checked.
  \item \texttt{[NEW]}: companion artifact produced in this project
    (Lean source, structural argument, or verification log). Listed
    separately from external prior art.
\end{itemize}
For the load-bearing analytic-NT citations, the current status is:
\textbf{Halberstam-Richert 1974}, \textbf{Iwaniec-Kowalski 2004},
\textbf{Pappalardi 2003}, and \textbf{Heath-Brown 1986} are at
\texttt{[EXISTS-LOCATION]}: paper + theorem number + the statement
shape are known and used; verbatim line-checking of the printed text
is a pending preprint task before arXiv submission. The four
``Action needed'' lines in L1--L4 above each name a specific
verbatim cross-check item.

\textbf{Action needed:} verbatim line-check of the four load-bearing
\texttt{[EXISTS-LOCATION]} citations (Pappalardi~2003, Halberstam-Richert
\S6 Thm~6.1, Iwaniec-Kowalski Thm~11.7, Iwaniec-Kowalski Ch.~22).

\paragraph{L8 --- The empirical-rigorous gap.}
The gap from rigorous $D_{0} \approx 10^{87}$ to empirical maximum
diagonal $D = 26$ (sweep complete to $m \leq 10^{10}$) is not
addressed. We conjecture but do not prove that the rigorous $D_{0}$
can be reduced to $D_{0} = O(\log m)$ in line with the empirical
growth.

\textbf{Action needed:} (optional, for a stronger result) close the
gap by exploiting the specific structure of $\langle 2, 3 \rangle$
beyond what we use here.

\paragraph{L9 --- Preprint status.}
This is a preprint, not a refereed publication. The author is
responsible for the content; literature citations are tagged in L7
above with their verification status. Source files for the working
drafts of each section are preserved alongside this preprint in the
source directory.

%==================== BIBLIOGRAPHY ====================

\begin{thebibliography}{99}

\bibitem[Bombieri 1965]{Bombieri1965}
E.~Bombieri.
\emph{On the large sieve.}
Mathematika 12 (1965), 201--225.
\texttt{[VERIFIED at paper existence and abstract]}

\bibitem[Bourgain-Glibichuk-Konyagin 2006]{BourgainGlibichukKonyagin2006}
J.~Bourgain, A.~Glibichuk, S.~V.~Konyagin.
\emph{Estimates for the number of sums and products and for
exponential sums in fields of prime order.}
J. London Math. Soc. (2) 73 (2006), no.~2, 380--398.
\texttt{[EXISTS-LOCATION]}

\bibitem[Brun 1915]{Brun1915}
V.~Brun.
\emph{\"Uber das Goldbachsche Gesetz und die Anzahl der Primzahlpaare.}
Arch. f. Math. og Naturvid. B34 (1915), no.~8, 1--19.
\texttt{[EXISTS-LOCATION; classical reference]}

\bibitem[Drmota-Tichy 1997]{DrmotaTichy1997}
M.~Drmota, R.~F.~Tichy.
\emph{Sequences, Discrepancies and Applications.}
Lecture Notes in Math. 1651, Springer, 1997.
\texttt{[EXISTS-LOCATION; standard reference for ET inequality]}

\bibitem[Erd\H{o}s-Graham 1980]{ErdosGraham1980}
P.~Erd\H{o}s, R.~L.~Graham.
\emph{Old and new problems and results in combinatorial number theory.}
L'Enseignement Math\'ematique Monographies 28, Geneva, 1980.
Problem~203, p.~27.
\texttt{[VERIFIED]}

\bibitem[Erd\H{o}s-Pomerance 1985]{ErdosPomerance1985}
P.~Erd\H{o}s, C.~Pomerance.
\emph{On the normal number of prime factors of $\phi(n)$.}
Rocky Mountain J. Math. 15 (1985), 343--352.
\texttt{[VERIFIED at paper existence; the bound we use is from the
Pomerance 1985 companion paper "On the distribution of multiplicative
orders"]}

\bibitem[Friedlander-Iwaniec 1998]{FriedlanderIwaniec1998}
J.~Friedlander, H.~Iwaniec.
\emph{The polynomial $X^{2} + Y^{4}$ captures its primes.}
Ann. of Math. (2) 148 (1998), no.~3, 945--1040.
\texttt{[VERIFIED at title and abstract; key prior art for
asymptotic-formula sieves at $\kappa = 1/2$]}

\bibitem[Greaves 2001]{Greaves2001}
G.~Greaves.
\emph{Sieves in number theory.}
Ergebnisse der Mathematik 43, Springer, 2001.
\texttt{[EXISTS-LOCATION; standard reference for sieve functions at
various $\kappa$]}

\bibitem[Halberstam-Richert 1974]{HalberstamRichert1974}
H.~Halberstam, H.-E.~Richert.
\emph{Sieve methods.}
London Math. Soc. Monographs 4, Academic Press, 1974.
\texttt{[EXISTS-LOCATION; standard reference]}

\bibitem[Heath-Brown 1986]{HeathBrown1986}
D.~R.~Heath-Brown.
\emph{Artin's conjecture for primitive roots.}
Quart. J. Math. Oxford Ser. (2) 37 (1986), no.~145, 27--38.
\texttt{[EXISTS-LOCATION; verified at paper existence + statement of
positive-density result]}

\bibitem[Heath-Brown 2001]{HeathBrown2001}
D.~R.~Heath-Brown.
\emph{Primes represented by $x^{3} + 2 y^{3}$.}
Acta Math. 186 (2001), no.~1, 1--84.
\texttt{[EXISTS-LOCATION]}

\bibitem[Heath-Brown-Konyagin 2000]{HeathBrownKonyagin2000}
D.~R.~Heath-Brown, S.~V.~Konyagin.
\emph{New bounds for Gauss sums derived from kth powers, and for
Heilbronn's exponential sum.}
Quart. J. Math. Oxford 51 (2000), 221--235.
\texttt{[EXISTS-LOCATION]}

\bibitem[Hooley 1972]{Hooley1972}
C.~Hooley.
\emph{On the Brun-Titchmarsh theorem.}
J. Reine Angew. Math. 255 (1972), 60--79.
\texttt{[EXISTS-LOCATION; for Brun-Hooley sieve improvement]}

\bibitem[Iwaniec 1980]{Iwaniec1980}
H.~Iwaniec.
\emph{Rosser's sieve.}
Acta Arith. 36 (1980), 171--202.
\texttt{[VERIFIED at paper existence + general framework; the $\kappa
= 0$ specialization we cite is not the explicit subject of this
paper]}

\bibitem[Iwaniec-Kowalski 2004]{IwaniecKowalski2004}
H.~Iwaniec, E.~Kowalski.
\emph{Analytic Number Theory.}
AMS Colloquium Publications 53, 2004.
\texttt{[EXISTS-LOCATION; Theorems 6.1 (Brun-Hooley), 11.6
(Rosser-Iwaniec linear sieve), 11.7 (large sieve), 22.x (asymptotic
sieve)]}

\bibitem[Maynard 2019]{Maynard2019}
J.~Maynard.
\emph{Primes with restricted digits.}
Invent. Math. 217 (2019), no.~1, 127--218.
\texttt{[VERIFIED at title and abstract]}

\bibitem[Mauduit-Rivat 2010]{MauduitRivat2010}
C.~Mauduit, J.~Rivat.
\emph{Sur un probl\`eme de Gelfond: la somme des chiffres des nombres
premiers.}
Ann. of Math. (2) 171 (2010), no.~3, 1591--1646.
\texttt{[VERIFIED at title and abstract]}

\bibitem[Pappalardi 1995]{Pappalardi1995}
F.~Pappalardi.
\emph{On the order of finitely generated subgroups of
$\mathbb{Q}^{\times}$ (mod $p$) and divisors of $p - 1$.}
J. Number Theory 57 (1996), no.~2, 207--222.
[Submitted 1995.]
\texttt{[EXISTS-LOCATION; Theorem 1 statement read at abstract level]}

\bibitem[Pappalardi 2003]{Pappalardi2003}
F.~Pappalardi.
\emph{Square-free values of the order function.}
J. Number Theory 100 (2003), 251--263.
\texttt{[EXISTS-LOCATION; rank-2 quantitative refinement]}

\bibitem[Pappalardi-Susa 2010]{PappalardiSusa2010}
F.~Pappalardi, A.~Susa.
\emph{An analogue of Artin's conjecture for multiplicative subgroups
of the rationals.}
Arch. Math. 95 (2010), no.~1, 27--37.
\texttt{[EXISTS-LOCATION]}

\bibitem[Pomerance 1985]{Pomerance1985}
C.~Pomerance.
\emph{On the distribution of multiplicative orders.}
1985 (informal note / appendix to Erd\H{o}s-Pomerance 1985).
\texttt{[EXISTS-LOCATION; the bound $\#\{p \leq z : \mathrm{ord}_{p}(a)
< z^{1-\alpha}\} \ll z^{1 - \alpha/2}$ is the standard
Erd\H{o}s-Pomerance bound]}

\bibitem[Rosser-Schoenfeld 1962]{RosserSchoenfeld1962}
J.~B.~Rosser, L.~Schoenfeld.
\emph{Approximate formulas for some functions of prime numbers.}
Illinois J. Math. 6 (1962), 64--94.
\texttt{[VERIFIED at paper and Theorem 23 statement; used for Mertens
explicit constants]}

\bibitem[Wilder 2026a]{WilderChain103}
J.~Wilder.
\emph{Exact structural characterization of uncovered cells in the
$(2520, 5040)$ period base for the V-family.}
Companion note to the present preprint, 2026. Includes a Lean 4
formalization of the universal density statement
(\texttt{Chain103UniversalDensity.lean}) and Python verification
across three independent period extensions.
\texttt{[NEW; Lean-formalized; awaiting peer review]}

\bibitem[Wilder 2026b]{WilderEG203Lean2026}
J.~Wilder.
\emph{Erd\H{o}s-Graham~\#203: a kernel-verified closure in Lean 4.}
Companion Lean 4 development, 2026.
\texttt{[NEW; Lean-formalized; awaiting peer review]}

\bibitem[Wilder 2026c]{WilderEG203MasterLog}
J.~Wilder.
\emph{Erd\H{o}s-Graham~\#203: a verification log.}
Companion technical note recording the empirical sweeps and
verification chains referenced in \S\ref{sec:empirical}, 2026.
\texttt{[NEW]}

\bibitem[Vinogradov 1965]{Vinogradov1965}
A.~I.~Vinogradov.
\emph{The density hypothesis for Dirichlet $L$-series.}
Izv. Akad. Nauk SSSR Ser. Mat. 29 (1965), 903--934.
\texttt{[EXISTS-LOCATION; classical BV co-discoverer reference]}

\end{thebibliography}

\end{document}
